\chapter{\label{chap:chap6}Considerações Finais}

Este volume consolidou o desenvolvimento, a validação e a implantação do \textbf{Dashboard Financeiro}, uma aplicação prática e funcional voltada à organização das finanças pessoais. O produto final encontra-se publicado em nuvem, oferece dashboards completos, módulo educacional e notificações automáticas, além de servir como laboratório para aplicar na prática os conhecimentos adquiridos ao longo da graduação em Sistemas de Informação.

O projeto representa não apenas uma entrega técnica, mas também uma jornada de crescimento pessoal e profissional. A construção dos microserviços, do frontend e do pipeline de deploy com domínio próprio demonstrou domínio das disciplinas de arquitetura, engenharia de dados e experiência do usuário, reforçando o aprendizado acumulado durante o curso.

\section*{Atendimento aos apontamentos do TC I}

A banca responsável pela avaliação do TC~I registrou observações críticas na etapa anterior; todas foram tratadas neste volume:
\begin{itemize}
    \item \textbf{Análise de mercado aprofundada (Seção 2.7).} O Capítulo~\ref{chap:chap2} foi ampliado com descrições detalhadas dos concorrentes, tabela comparativa e subseção específica mostrando como o Dashboard Financeiro cobre as lacunas identificadas.
    \item \textbf{Estrutura textual consistente.} A organização dos capítulos foi revisada, eliminando repetições (objetivos, funcionalidades e arquitetura aparecem uma única vez nos Capítulos~\ref{chap:intro}--\ref{chap:chap3}). O Capítulo~\ref{chap:chap4} passou a conter apenas cronograma, rastreabilidade e validação das funcionalidades.
    \item \textbf{Modelagem semântica.} A Figura~\ref{fig:arquitetura-geral} segue a notação C4 (nível \textit{Container}); a Figura~\ref{fig:modelo-dominio} apresenta o modelo de domínio antes do banco; a Figura~\ref{fig:casos-uso} utiliza UML formal e está compatibilizada com a lista de funcionalidades e descrições tipo história de usuário.
    \item \textbf{Casos de uso x funcionalidades}. A Tabela~\ref{tab:rastreabilidade-funcionalidades} (Capítulo~\ref{chap:chap4}) vincula cada funcionalidade planejada aos módulos entregues, e o diagrama de casos de uso reflete exatamente esse escopo.
    \item \textbf{Definição do MVP.} O Capítulo~\ref{chap:chap3} ganhou a seção ``Definição do MVP'', descrevendo os itens construídos entre junho e agosto (cadastro, contas, transações, metas, dashboards e educação).
    \item \textbf{Entrega em nuvem e viabilidade.} O Capítulo~\ref{chap:chap5} descreve o deploy real na Railway, o domínio público, a autenticação de DNS na KingHost e a operação do serviço de notificações via Resend, evidenciando a capacidade de atingir os objetivos propostos.
\end{itemize}

\section*{Impacto Esperado}

Com a aplicação publicada no domínio dedicado do projeto, os usuários já dispõem de uma ferramenta que sintetiza dados financeiros pessoais em visões compreensíveis, enviando alertas contextualizados diretamente para o e-mail verificado. Segundo o Banco Central do Brasil, o acesso à informação financeira é essencial para a autonomia econômica e a redução do endividamento, conforme o Relatório de Cidadania Financeira 2023 \cite{BCB2023CIDADANIA}.

Os achados do Cetic.br reforçam que a compreensão dos próprios hábitos, aliada a ferramentas digitais, reduz o endividamento crônico \cite{CETIC2023}. O módulo educacional e os dashboards do sistema permitem que o usuário identifique padrões, corrija excessos e trace estratégias conscientes, materializando o impacto discutido no Capítulo~\ref{chap:chap2}.

\section*{Contribuição Social e Educacional}

Mais do que um software, o Dashboard Financeiro atua como ferramenta de transformação. Relatórios internacionais, como o estudo da OCDE sobre alfabetização financeira, apontam que programas de educação financeira reduzem desigualdades e ampliam o bem-estar coletivo \cite{OECD2023}. Ao integrar orientação contextual (módulo \textit{Educação}) e visualizações acionáveis, o sistema catalisa mudanças comportamentais com base em dados reais.

Segundo Gitman e Joehnk, o planejamento financeiro pessoal é essencial para alcançar objetivos de longo prazo e garantir estabilidade econômica \cite{GITMAN2018}. Alinhado a essa visão, o Dashboard Financeiro auxilia usuários a definir metas tangíveis, acompanhar indicadores e celebrar pequenas vitórias, reforçando a formação continuada em educação financeira.

O projeto também demonstra o compromisso acadêmico com tecnologias acessíveis e impacto social: o código é versionado publicamente, o domínio próprio democratiza o acesso e o pipeline de notificações amplia o alcance das recomendações.

\section*{Limitações Encontradas}

Durante o desenvolvimento, destacaram-se as seguintes limitações:
\begin{itemize}
    \item A integração com instituições financeiras permanece manual (entrada de dados pelo usuário), uma vez que o escopo do TC~II não contemplou integrações Open Finance;
    \item O monitoramento em produção depende de logs da Railway e da Kafka UI local; seria interessante incorporar métricas centralizadas (Prometheus/Grafana) em uma próxima etapa;
    \item O módulo educacional fornece conteúdo textual, mas não há acompanhamento estatístico do engajamento ou experimentos controlados com usuários reais.
\end{itemize}

Mesmo com tais restrições, o produto entregue é funcional e estabelece uma base sólida para evoluções.

\section*{Desafios Enfrentados}

Alguns desafios técnicos marcaram esta jornada:
\begin{itemize}
    \item Integrar o frontend ao API Gateway e este aos microserviços via HTTP, enquanto os serviços trocavam eventos pelo Kafka, exigiu ajustes finos em autenticação, roteamento e contratos assíncronos;
    \item A configuração do domínio e dos registros DNS na KingHost demandou múltiplas validações (SPF, DKIM, CNAME), além de sincronizar essas informações com o gateway e com a plataforma de e-mails;
    \item O deploy no Railway com imagens Docker foi trabalhoso porque o projeto depende de diversos componentes de infraestrutura (Kafka, Zookeeper, MySQL), impondo muitas variáveis de ambiente e regras de rede interna dentro da plataforma.
\end{itemize}

\section*{Trabalhos Futuros}

Como perspectivas de continuidade e evolução, propõem-se:
\begin{itemize}
    \item Integração com APIs bancárias (Open Finance) para importação automática e segura de transações;
    \item Inclusão de inteligência artificial para análise preditiva do comportamento financeiro e recomendações personalizadas;
    \item Desenvolvimento de uma versão mobile nativa com distribuição nas lojas de aplicativos;
    \item Monitoramento contínuo de performance e escalabilidade com uso de Prometheus, Grafana e alertas pró-ativos em produção.
\end{itemize}

Esses aprimoramentos visam tornar o \textbf{Dashboard Financeiro} ainda mais completo, acessível e impactante, contribuindo efetivamente para o bem-estar financeiro da população.
